% A printable figure deck: exactly one blueprint illustration per A4 page.
% latexmkrc builds this alongside print.tex as blueprint/print/illustrations.pdf.
\documentclass[a4paper]{article}

\usepackage[margin=18mm]{geometry}
\usepackage{amssymb, amsmath}
\usepackage{graphicx}
\usepackage{tikz}
\usetikzlibrary{arrows.meta,decorations.markings}

\pagestyle{empty}
\setlength{\parindent}{0pt}
% Generated by `lake exe gluckFigures`; do not edit.
\def\EuclideanOpenA{0.25}
\def\EuclideanB{0.1666666667}
\def\SphereRadius{0.5}
\def\SphereCurvature{0.75}
\def\HyperbolicCOne{2}
\def\HyperbolicCTwo{1.2}
\def\HyperbolicCThree{1.02}


\newcommand{\illustrationpage}[3]{%
  \begin{center}
    \vspace*{\fill}
    {\LARGE\bfseries #2\par}
    \vspace{1.25cm}
    \resizebox{0.98\textwidth}{!}{\input{figures/#1.tikz}}%
    \vspace{1.1cm}

    \begin{minipage}{0.86\textwidth}
      \centering\large #3
    \end{minipage}
    \vspace*{\fill}
  \end{center}%
}

\begin{document}

\illustrationpage{euclidean-closure}
  {Euclidean closure defect}
  {The first Fourier mode of the radius-of-curvature weight is the
   translational closure defect. Removing it changes the displayed
   reconstruction from open to closed.}

\newpage
\illustrationpage{euclidean-gluck-reconstructions}
  {Three Euclidean Gluck reconstructions}
  {Three positive curvature functions and the closed convex curves reconstructed
   from their reciprocal radius-of-curvature weights. With no first Fourier
   harmonic, each displayed reconstruction closes. Matching numbered markers
   identify four selected curvature extrema and their corresponding curve vertices.}

\newpage
\illustrationpage{euclidean-dahlberg-reconstructions}
  {Three Euclidean Dahlberg reconstructions}
  {Three mixed-sign curvature functions and their unit-speed reconstructed
   curves. Matching numbered markers identify two positive local maxima and two
   local minima; additional unmarked vertices are allowed. Half-turn symmetry
   makes each displayed reconstruction close.}

\newpage
\illustrationpage{sphere-stereographic}
  {Spherical stereographic curvature}
  {The conformal-curvature law and its centered-circle sign check at
   $r=1/2$, where $\kappa_S=3/4$ and the tangent-angle gauge speed is $q=r$.}

\newpage
\illustrationpage{sphere-reconstructions}
  {Spherical positive and mixed reconstructions}
  {Two positive spherical geodesic-curvature profiles and one mixed-sign profile,
   paired with their inverse-stereographic lifts on the unit sphere. Four selected vertices are
   matched in each example; additional extrema remain unmarked.}

\newpage
\illustrationpage{hyperbolic-escape}
  {Hyperbolic escape threshold}
  {Centered hyperbolic circles approach the ideal boundary as
   $c\downarrow1$, while the admissibility bracket
   $\sqrt{c^2-1}$ collapses to zero.}

\newpage
\illustrationpage{hyperbolic-reconstructions}
  {Hyperbolic positive and mixed reconstructions}
  {Two profiles above the hyperbolic escape threshold $\kappa_H=1$ and one
   mixed-sign profile, paired with their Poincaré-disk representatives. Four
   selected vertices are matched in each example.}

\newpage
\illustrationpage{winding-angle-lift}
  {Winding number and the lifted angle}
  {A nonvanishing complex loop is normalized to the circle and lifted to a
   real angle. The lift's endpoint displacement records the integer winding
   number; this example winds twice.}

\newpage
\illustrationpage{winding-error-map}
  {Boundary winding forces closure}
  {The configuration disk maps to endpoint-error space. Nonzero winding of
   the boundary error loop forces an interior zero, hence a parameter whose
   reconstructed curve has zero endpoint defect.}

\newpage
\illustrationpage{degree-one-reparametrization}
  {Degree-one reparametrization}
  {An increasing periodic lift redistributes parameter time while preserving
   cyclic order, curvature values, and the correspondence between vertices.}

\newpage
\illustrationpage{hyperbolic-geometry-atlas}
  {Hyperbolic metric and constant-curvature atlas}
  {The Poincaré metric amplifies Euclidean motion near the ideal boundary.
   Geodesics, equidistant curves, horocycles, and circles occupy the four
   constant-geodesic-curvature regimes separated by the threshold one.}

\newpage
\illustrationpage{hyperbolic-flow-stability}
  {Hyperbolic confinement and perturbation stability}
  {An inward-pointing field preserves the admissible domain. An $L^1$ profile
   bound controls trajectory displacement, allowing even deep negative spikes
   when their integrated effect is sufficiently small.}

\newpage
\illustrationpage{hyperbolic-closing-engine}
  {The hyperbolic closing engine}
  {Mirror-axis quarter landing, Poincaré--Miranda face signs, five-leg
   transport, and a unique turning root combine to close the coupled
   hyperbolic reconstruction.}

\newpage
\illustrationpage{hyperbolic-simplicity}
  {Chord nonvanishing and simplicity}
  {Every proper sub-arc has a nonzero integrated tangent chord. A
   self-intersection would produce a zero chord between two distinct parameter
   times and is therefore excluded.}

\newpage
\illustrationpage{spaceform-unification}
  {One construction across all three space forms}
  {A single profile and its marked vertices are followed through spherical,
   Euclidean, and hyperbolic reconstruction. The curvature law and geometry
   vary continuously with ambient curvature $K$ through the flat case.}

\end{document}
